\documentclass[a4paper,12pt]{article}

\usepackage[utf8]{inputenc}
\usepackage[spanish, es-noshorthands]{babel}
\usepackage{geometry}
\usepackage{graphicx}
\usepackage{multirow}
\usepackage{array}
\usepackage{fancyhdr}
\usepackage{lastpage}
\usepackage{hyperref}
\usepackage{float}
\usepackage{caption}
\usepackage{booktabs}
\usepackage[table]{xcolor}
\usepackage[T1]{fontenc}
\usepackage{tabularx}

\hypersetup{colorlinks=true,linkcolor=blue,urlcolor=blue}

\definecolor{SKTPrimary}{HTML}{3F4A4F}
\definecolor{SKTDark}{HTML}{242B2E}
\definecolor{SKTAccent}{HTML}{979C9D}
\definecolor{SKTLight}{HTML}{CDCFD2}
\definecolor{SKTMid}{HTML}{6E7375}

\geometry{
top=2.5cm,
bottom=2.5cm,
left=3cm,
right=3cm,
includehead,
headheight=130pt,
headsep=1cm
}

% ---------------- ENCABEZADO ----------------
\pagestyle{fancy}
\fancyhf{}
\renewcommand{\headrulewidth}{0pt}

\fancyhead[C]{%
\small
\setlength{\tabcolsep}{4pt}
\renewcommand{\arraystretch}{1.1}
\begin{tabular}{|m{3cm}|m{7.5cm}|m{3.8cm}|}
\hline
\multirow{5}{*}{
\parbox[c][3.5cm][c]{2.5cm}{
\centering
\includegraphics[width=2.2cm]{SKTLogo.png}
}
}
& \centering\textbf{SKT Software Solution}
& \parbox[c]{\linewidth}{\centering
\textbf{CÓDIGO:}\\
MI-DES-PLA-004
\vspace{3pt}
} \\ \cline{2-3}

& \centering\textbf{PROCESO: Desarrollo de Software}
& \parbox[c]{\linewidth}{\centering
\textbf{VERSIÓN:}\\
1.0
\vspace{3pt}
} \\ \cline{2-3}

& \centering Planificación Sprint 3
& \parbox[c]{\linewidth}{\centering
\textbf{ELABORACIÓN:}\\
11/04/2026
\vspace{3pt}
} \\ \cline{2-3}

& \centering Motor de Agendamiento e Interfaces Frontend
& \parbox[c]{\linewidth}{\centering
\textbf{ÚLTIMA REVISIÓN:}\\
11/04/2026
} \\ \hline
\end{tabular}%
}

\begin{document}
\raggedbottom

\section{Título}
Planificación del Sprint 3: Motor de Agendamiento, Concurrencia e Interfaces Frontend - LUPSI

\section{Introducción}
El presente documento formaliza la planificación del tercer sprint del proyecto LUPSI. Su propósito es dejar claramente definidos el objetivo del sprint, el alcance técnico, las actividades prioritarias, los entregables esperados y su alineación con el cronograma general.

Tras haber consolidado la base de datos, la seguridad (RLS) y la autenticación en el Sprint 2, el Sprint 3 aborda el núcleo operativo del negocio: la gestión de citas médicas. Esta etapa materializa las historias de usuario US-02 y US-03, garantizando un agendamiento síncrono sin conflictos y proveyendo las interfaces gráficas (PWA) tanto para los pacientes como para el personal administrativo.

\section{Objetivo}
Desarrollar e integrar el Motor de Reservas en el backend para prevenir el cruce de horarios, y construir los portales web progresivos (Frontend) que permitan a los pacientes agendar citas y al personal de recepción visualizar la carga diaria de forma segura y en tiempo real.

\section{Alcance del Sprint 3}
El Sprint 3 se centra exclusivamente en la lógica de negocio de agendamiento y en las vistas operativas de los usuarios.

\subsection{Incluye}
\begin{itemize}
    \item Lógica backend de agendamiento síncrono (Motor de reservas).
    \item Control de concurrencia en base de datos para prevención de cruce de horarios.
    \item Desarrollo Frontend (PWA) del Portal del Paciente (Agendamiento e historial).
    \item Desarrollo Frontend del Panel Administrativo (Vista de agenda diaria para recepción).
    \item Integración del módulo de autenticación JWT del Sprint 2 con las vistas Frontend.
\end{itemize}

\subsection{No incluye (Asignado al Sprint 4)}
\begin{itemize}
    \item Módulo de Historia Clínica Digital.
    \item Integración con Cloudinary para almacenamiento de recetas médicas (PDFs).
    \item Manejo de indicaciones post-consulta (RBAC específico).
\end{itemize}

\section{Definiciones}
\begin{itemize}
    \item \textbf{Sprint:} periodo de trabajo definido dentro de Scrum para entregar un conjunto acotado de tareas.
    \item \textbf{PWA (Progressive Web App):} Aplicación web que ofrece una experiencia similar a una app nativa en dispositivos móviles.
    \item \textbf{Concurrencia:} Capacidad del sistema para manejar múltiples solicitudes de agendamiento al mismo tiempo sin colisionar datos.
    \item \textbf{JWT (JSON Web Token):} Estándar para la transmisión segura de información de autenticación entre el cliente y el servidor.
\end{itemize}

\section{Planificación del Sprint 3}

\subsection{Objetivo específico del sprint}
Implementar el ciclo completo de agendamiento, asegurando que un paciente pueda reservar una cita desde la interfaz visual y el backend la valide y persista sin permitir que dos personas ocupen el mismo horario.

\subsection{Entregables del Sprint}

\footnotesize
\renewcommand{\arraystretch}{1.2}
\begin{tabularx}{\textwidth}{|>{\raggedright\arraybackslash}X|>{\centering\arraybackslash}p{2.5cm}|>{\centering\arraybackslash}p{2.0cm}|>{\centering\arraybackslash}p{1.5cm}|}
\hline
\textbf{Entregable} & \textbf{Responsable} & \textbf{Fecha} & \textbf{Aceptado} \\
\hline
1. Backend: Motor de Reservas y control de Concurrencia & Sebastián Ortiz & 21/04/2026 &  \\
\hline
2. Frontend: Portal del Paciente (PWA) & Daniel Luisa & 30/04/2026 &  \\
\hline
3. Frontend: Panel Administrativo & Alex Guachi & 30/04/2026 &  \\
\hline
4. Integración Auth-Frontend (Deuda Técnica S2) & Alex Guachi & 15/04/2026 &  \\
\hline
\end{tabularx}
\normalsize

\subsection{Backlog del Sprint 3 (Desglose de Tareas)}
La siguiente tabla detalla las tareas operativas específicas (T-XX) a ejecutar durante este sprint y su trazabilidad directa con los requerimientos aprobados en el Product Backlog (PB-XX).

\renewcommand{\arraystretch}{1.25}
\begin{tabularx}{\textwidth}{|p{0.08\textwidth}|p{0.12\textwidth}|X|p{0.18\textwidth}|p{0.15\textwidth}|}
\hline
\textbf{Tarea} & \textbf{Origen} & \textbf{Actividad (Sprint Backlog)} & \textbf{Responsable} & \textbf{Estado} \\
\hline
T-01 & PB-11 & Desarrollar lógica de agendamiento síncrono previniendo cruce de horarios. & Sebastián & En Curso \\
\hline
T-02 & PB-12 & Construir interfaces PWA para reserva y consulta de citas del paciente. & Daniel & En Curso \\
\hline
T-03 & PB-13 & Crear vista de agenda diaria y administración para personal de recepción. & Alex & En Curso \\
\hline
T-04 & Deuda S2 & Integrar las interfaces web del login con la lógica del servidor API/JWT. & Alex & En Curso \\
\hline
\end{tabularx}

\subsection{Cronograma Detallado}
\small
\renewcommand{\arraystretch}{1.2}
\begin{tabularx}{\textwidth}{|>{\centering\arraybackslash}p{2.2cm}|>{\centering\arraybackslash}p{0.7cm}|>{\centering\arraybackslash}p{0.9cm}|>{\centering\arraybackslash}p{0.9cm}|>{\centering\arraybackslash}p{0.7cm}|>{\centering\arraybackslash}p{1.2cm}|X|X|}
\hline
\textbf{Tarea} & \textbf{Días} & \textbf{Inicio} & \textbf{Fin} & \textbf{H} & \textbf{Hito} & \textbf{Entregable} & \textbf{Responsable} \\
\hline
Integración Auth & 6 & 08/04 & 15/04 & 24 & Hito 1 & Login PWA & Alex \\
\hline
Motor Reservas & 10 & 08/04 & 21/04 & 40 & Hito 1 & API Agendamiento & Sebastián \\
\hline
Portal Paciente & 16 & 08/04 & 30/04 & 64 & Hito 2 & Vistas Frontend & Daniel \\
\hline
Panel Admin & 16 & 08/04 & 30/04 & 64 & Hito 2 & Vistas Recepción & Alex \\
\hline
\end{tabularx}
\normalsize

\section{Criterios de aceptación del Sprint 3}
El sprint se considerará aprobado cuando se cumpla lo siguiente:
\begin{itemize}
    \item El sistema backend rechaza estrictamente la superposición de citas médicas (US-02).
    \item El Portal del Paciente permite visualizar horarios disponibles y confirmar reservas de forma exitosa.
    \item La agenda administrativa refleja en tiempo de ejecución los horarios reservados sin exponer data clínica sensible (US-03).
    \item Las interfaces frontend están conectadas de manera segura con el backend mediante los tokens JWT.
\end{itemize}

\section{Trazabilidad con las historias de usuario}
Este sprint entrega valor directo y funcional enfocado en el core del negocio:
\begin{itemize}
    \item \textbf{US-02 (Agendamiento Síncrono):} Implementado a través del Motor de Reservas (PB-11) y el Portal del Paciente (PB-12).
    \item \textbf{US-03 (Visualización de Agenda Administrativa):} Implementado a través del Panel Administrativo web (PB-13).
\end{itemize}

\section{Riesgos del Sprint 3}
\begin{itemize}
    \item Retrasos en el desarrollo del Motor de Reservas (PB-11) bloquearán las pruebas de integración del Frontend de Daniel y Angel.
    \item Errores en la gestión de transacciones en la base de datos que permitan agendar dos pacientes a la misma hora (condición de carrera).
    \item Problemas de configuración CORS (Cross-Origin Resource Sharing) al comunicar los repositorios locales de Frontend y Backend.
\end{itemize}

\section{Criterios generales de calidad}
Para considerar el sprint como correctamente ejecutado, la solución debe cumplir con los siguientes criterios:
\begin{itemize}
    \item Tiempos de respuesta óptimos en las consultas de disponibilidad de horarios.
    \item Retroalimentación visual clara en el frontend cuando un paciente reserva exitosamente o si ocurre un error.
    \item Código fuente documentado y subido a las ramas principales correspondientes en GitHub.
\end{itemize}

\section{Conclusiones}
El Sprint 3 constituye el hito más importante para la viabilidad operativa del proyecto LUPSI, ya que resuelve el problema central del centro médico: la gestión de citas sin conflictos. 

Al finalizar este sprint, el sistema dejará de ser una estructura base para convertirse en una aplicación web progresiva funcional, permitiendo la interacción real entre pacientes y administradores, y dejando el terreno preparado para la digitalización de documentos clínicos en el Sprint 4.

\section{Firmas de Responsabilidad}

\renewcommand{\arraystretch}{2}

\noindent
\begin{tabular}{
p{0.28\textwidth}
p{0.30\textwidth}
>{\centering\arraybackslash}p{0.22\textwidth}
>{\centering\arraybackslash}p{0.12\textwidth}
}
\hline
\textbf{Rol} &
\textbf{Nombre / Representante} &
\textbf{Firma (QR)} &
\textbf{Fecha} \\
\hline

Gestor del Proyecto
& Angel Ayuquina
& \includegraphics[width=3.2cm]{QR_Sello_Angel_Ayuquina.png}
& 11/04/2026 \\[0.6cm]

Desarrollador Backend
& Sebastián Ortiz
& \includegraphics[width=3.2cm]{QR_Sello_Sebastían_Ortiz.png}
& 11/04/2026 \\[0.6cm]

Desarrollador Fullstack
& Daniel Luisa
& \includegraphics[width=3.2cm]{QR_Sello_Daniel_Luisa.png}
& 11/04/2026 \\[0.6cm]

Desarrollador Frontend
& Alex Guachi
& \includegraphics[width=3.2cm]{QR_Sello_Alex_Huachi.png}
& 11/04/2026 \\[0.6cm]

\hline
\end{tabular}

\section{Control de cambios}

\small
\renewcommand{\arraystretch}{1.25}
\begin{tabularx}{\textwidth}{|p{0.12\textwidth}|X|p{0.25\textwidth}|p{0.18\textwidth}|}
\hline
\textbf{Versión} & \textbf{Descripción del Cambio} & \textbf{Responsable del Cambio} & \textbf{Fecha de Aprobación} \\
\hline
1.0 & Elaboración de la planificación específica para el Sprint 3 (Motor de Reservas y Frontend). & Angel Ayuquina & 11/04/2026 \\
\hline
\end{tabularx}
\normalsize

\end{document}